\documentclass[10pt,a4paper,openany,twoside, eleve]{book}

\usepackage{layout_cours}
\setlist{leftmargin=5mm}

\begin{document}
\chapnumber{} % numéro du chapitre
\titre{Calcul avec les racines carrées} % titre du chapitre
\classe{Seconde} % classe
\entetegal

% Debut du texte
% **************

\setcounter{section}{1}
\section*{Exercices}
\begin{exercicesnofoot}
	\sectionexo{Echauffement}
	\begin{exo}
		Calculer les expressions suivantes :
			\begin{enumerate}[label={}]
				\item $A = \sqrt{3} \times 4\sqrt{3}=$
				\item $B = \left(2\sqrt{7}\right)\times\left(-5\sqrt{7}\right)=$
				\item $C = \left(3\sqrt{7}\right)^2=$
				\item $D = \left(-2\sqrt{5}\right)\times\left(-4\sqrt{5}\right)=$
				\item $E = \sqrt{4\times 9\times 16}=$
				\item $F = \sqrt{2^2\times 3^4\times 25}=$
				\item $G = \sqrt{144\times 9}=$
				\item $H = \sqrt{\frac{225\times 121}{36 \times 49}}=$
			\end{enumerate}
	\end{exo}

	\sectionexo{Simplifier une racine carrée}
	\begin{exo}
		\begin{multicols}{2}
		\begin{enumerate}[label={}]
			\item $1^2 = \troum{1}$
			\item $2^2 = \troum{4}$
			\item $3^2 = \troum{9}$
			\item $4^2 = \troum{16}$
			\item $5^2 = \troum{25}$
			\item $6^2 = \troum{36}$
			\item $7^2 = \troum{49}$
			\item $8^2 = \troum{64}$
			\item $9^2 = \troum{81}$
			\item $10^2 = \troum{100}$
			\item $11^2 = \troum{121}$
			\item $12^2 = \troum{144}$
			\item $13^2 = \troum{169}$
			\item $14^2 = \troum{196}$
			\item $15^2 = \troum{225}$
			\item $16^2 = \troum{256}$
			\item $17^2 = \troum{289}$
			\item $18^2 = \troum{324}$
			\item $19^2 = \troum{361}$
			\item $20^2 = \troum{400}$
			\end{enumerate}
		\end{multicols}
	\end{exo}

	\begin{exo}
		\begin{multicols}{2}
			\begin{enumerate}[label={}]
			\item $18 = \troum{3^2 \times 2}$
			\item $12 = \troum{2^2 \times 3}$
			\item $24 = \troum{2^2 \times 6}$
			\item $28 = \troum{2^2 \times 7}$
			\item $45 = \troum{3^2 \times 5}$
			\item $72 = \troum{6^2 \times 2}$
			\item $150 = \troum{5^2 \times 6}$
			\item $675 = \troum{15^2 \times 3}$
			\item $288 = \troum{12^2 \times 2}$
			\item $588 = \troum{14^2 \times 3}$
			\end{enumerate}
			\end{multicols}
				\end{exo}

	\columnbreak
	\begin{exo}
		\begin{enumerate}
			\item Écrire sous la forme «$a\sqrt{3}$»
		\begin{multicols}{2}
			\begin{enumerate}[label={}]
				\item $\sqrt{12} = \troum{2\sqrt{3}}$
				\item $\sqrt{27} = \troum{3\sqrt{3}}$
				\item $\sqrt{300} = \troum{10\sqrt{3}}$
				\item $\sqrt{192} = \troum{8\sqrt{3}}$
				\end{enumerate}
		\end{multicols}
			\item Écrire sous la forme «$a\sqrt{6}$»
		\begin{multicols}{2}
			\begin{enumerate}[label={}]
					\item $\sqrt{96} = \troum{4\sqrt{6}}$
					\item $\sqrt{150} = \troum{5\sqrt{6}}$
					\item $\sqrt{216} = \troum{6\sqrt{6}}$
					\item $\sqrt{384} = \troum{8\sqrt{6}}$
			\end{enumerate}
		\end{multicols}
		\end{enumerate}
	\end{exo}
	\begin{exo}{ Simplifier les racines carrées suivantes. Autrement dit, écrire sous la forme $a\sqrt{b}$ avec $a$ et $b$ entiers, $b$ étant le plus petit possible
		\begin{multicols}{2}
			\begin{enumerate}[label={}]
					\item $\sqrt{40} = \troum{2\sqrt{10}}$
					\item $\sqrt{99} = \troum{3\sqrt{11}}$
					\item $\sqrt{54} = \troum{3\sqrt{6}}$
					\item $\sqrt{63} = \troum{3\sqrt{7}}$
					\item $\sqrt{32} = \troum{4\sqrt{2}}$
					\item $\sqrt{288} = \troum{12\sqrt{2}}$
					\item $\sqrt{675} = \troum{15\sqrt{3}}$
					\item $\sqrt{72} = \troum{6\sqrt{2}}$
					\item $\sqrt{845} = \troum{13\sqrt{5}}$
					\item $\sqrt{847} = \troum{11\sqrt{7}}$
			\end{enumerate}
		\end{multicols}
	}
	\end{exo}
	
	\begin{exo}
		Comparer, sans utiliser la calculatrice, les nombres suivants :
		\begin{multicols}{2}
			\begin{enumerate}[label=$\alph*)$]
				\item $6\sqrt{2}$ $\dots$ $5\sqrt{3}$ 
				\item $7\sqrt{2}$ $\dots$ $3\sqrt{11}$
				\item $-3\sqrt{7}$ $\dots$ $5\sqrt{2}$
				\item $4\sqrt{7}$ $\dots$ $7\sqrt{4}$
	\end{enumerate}
	\end{multicols}

	\end{exo}
	\sectionexo{Fractions}
	\begin{exo}
	Écrire sous la forme $\frac{a\sqrt{b}}{c}$ avec $a$, $b$ et $c$ entiers:
	\begin{multicols}{2}
		\begin{enumerate}[label={}]
			\item $\frac{1}{\sqrt{5}} = \troum{\frac{\sqrt{5}}{5}}$
			\item $\frac{2}{\sqrt{3}} = \troum{\frac{2\sqrt{3}}{3}}$
			\item $\frac{4}{\sqrt{12}} = \troum{\frac{2\sqrt{3}}{3}}$
			\item $\frac{\sqrt{2}}{\sqrt{3}} = \troum{\frac{\sqrt{6}}{3}}$
			\item $\frac{\sqrt{5}}{\sqrt{36}} = \troum{\frac{\sqrt{5}}{6}}$
			\item $\frac{\sqrt{22}}{\sqrt{11}} = \troum{\sqrt{2}}$
			\item $\sqrt{\frac{4}{5}} = \troum{\frac{2\sqrt{5}}{5}}$
			\item $\sqrt{\frac{7}{2}} = \troum{\frac{\sqrt{14}}{2}}$
			\item $\sqrt{\frac{9}{7}} = \troum{\frac{3\sqrt{7}}{7}}$
			\item $\frac{\sqrt{16}\sqrt{5}}{2\sqrt{3}} = \troum{\frac{2\sqrt{15}}{3}}$
		\end{enumerate}
	\end{multicols}
	\end{exo}

	\begin{defi}{Quantités conjuguées}{}
		Les nombres $x=\sqrt{a}-\sqrt{b}$ et $y=\sqrt{a}+\sqrt{b}$ sont appelés quantités conjugues ou simplement \textbf{conjugués}. 
		
		Ils permettent de simplifier les fractions car:

		\begin{align*}
			xy 
			&=  \troum{\left(\sqrt{a}-\sqrt{b}\right)\left(\sqrt{a}+\sqrt{b}\right)} \\
			&= \troum{a-b}
		\end{align*}
	\end{defi}

	\begin{exo}
		Écrire la fraction sans racine au dénominateur :
			\begin{enumerate}[label={}]
			\item $\frac{1}{\sqrt{3}+\sqrt{5}} = \troum{\frac{\sqrt{5}-\sqrt{3}}{2}}$
			\item $\frac{2}{\sqrt{7}-\sqrt{2}} = \troum{\frac{2(\sqrt{7}+\sqrt{2})}{5}}$
			\item $\frac{5}{\sqrt{5}-2\sqrt{2}} = \troum{\frac{-5}{3}(\sqrt{5}+2\sqrt{2})}$
			\item $\frac{\sqrt{2}+\sqrt{3}}{\sqrt{3}-\sqrt{2}} = \troum{5 + 2\sqrt{6}}$
			\item $\frac{\sqrt{5}}{\sqrt{8}-\sqrt{2}} = \troum{\frac{\sqrt{10}}{2}}$
			\item $\frac{\sqrt{2}-\sqrt{5}}{\sqrt{3}+\sqrt{5}} = \troum{\frac{\sqrt{10} + \sqrt{15} - \sqrt{6} - 5}{2}}$
			\end{enumerate}
		\end{exo}
	

	\sectionexo{Identité remarquable}
	\begin{exo}{
		Developper et simplifier les expressions suivantes :
			\begin{enumerate}[label={}]
				\item $(4\sqrt{5} + 2)^2 = \troum{84 + 16\sqrt{5}}$
				\item $(-\sqrt{9} - 3\sqrt{7})^2 = \troum{72 + 18\sqrt{7}}$
				\item $(3 - 7\sqrt{11})(3 + 7\sqrt{11}) = \troum{-530}$
				\item $(6 + \sqrt{12})^2 = \troum{48 + 24\sqrt{3}}$
			\end{enumerate}
		}
	\end{exo}

	\sectionexo{Pour aller plus loin}
	\begin{exo}
		Pour tout réel $x$ positif, déterminer l'inverse de 
		
		$\sqrt{x+1}-\sqrt{x}$. Cet inverse sera écrit sans fraction.
	\end{exo}

	\begin{exo}
		On considère deux nombres réels positifs $a=\frac{x+y}{2}$ et $g=\sqrt{xy}$.

		Comparer les nombres $a$ et $g$.

		\emph{Remarque : $a$ est appelée la moyenne arithmétique et $g$ la moyenne géométrique des deux réels $x$ et $y$.}
	\end{exo}

	% Exprimer sous la forme la plus simple possible, d'une seule fraction irr\'eductible, sans racine carrée au dénominateur, et les expressions algébriques développées:

	% \begin{enumerate}
	% 	\item $\dfrac{15}{\sqrt{5}}$
	% 	\item $\left(\sqrt{12}-\sqrt{3}\right)^2$
	% 	\item $(3\sqrt{2})^2-(\sqrt{2}-1)^2$
	% 	\item $\dfrac{2}{2+\sqrt{3}}$
	% 	\item $\dfrac{\sqrt{3}+\sqrt{12}}{\sqrt{3}-\sqrt{12}}$
	% \end{enumerate}

\end{exercicesnofoot}

\end{document}

